
\chapter{Sistemi lineari e interpolazione}
\newpage
\minitoc
\newpage

\section{Richiami teorici}
\label{sec:term-structure}

Un'altra applicazione dei metodi numerici per la soluzione dei
sistemi lineari riguarda la finanza e, in particolare, la
costruzione della \emph{struttura a termine dei tassi di interesse}.

Il problema che vogliamo affrontare può essere formulato nel modo
seguente: osservando sul mercato i prezzi di un insieme di
obbligazioni e conoscendo i flussi di cassa che esse genereranno
in futuro, possiamo ricavare i tassi di interesse che il mercato
associa alle diverse scadenze?

Come vedremo, anche questo problema può essere ricondotto alla
soluzione di un sistema di equazioni lineari. In questo modo,
lo stesso problema numerico
$$Ax=b$$
che abbiamo incontrato nelle applicazioni precedenti, ricompare in un contesto finanziario.


\subsection{Valore del denaro nel tempo e fattori di sconto}

Un principio fondamentale della finanza è che il valore di una
somma di denaro dipende dal momento nel quale essa diventa
disponibile. In generale, un euro disponibile oggi vale più di un
euro che verrà ricevuto in futuro.

Supponiamo di ricevere con certezza un ammontare $FV$ tra 
$T$ periodi. Se il tasso di interesse relativo a quella scadenza è $r_T$,
il valore attuale del pagamento futuro è

\begin{equation}
    PV=
    \frac{FV}{(1+r_T)^T}.
    \label{eq:present-value-term}
\end{equation}

Possiamo separare il pagamento futuro dal termine che permette
di riportarlo al suo valore attuale definendo il
\emph{fattore di sconto}

\begin{equation*}
    d_T=
    \frac{1}{(1+r_T)^T}.
    \label{eq:discount-factor}
\end{equation*}

L'Equazione~\eqref{eq:present-value-term} può quindi essere
riscritta semplicemente come

\begin{equation*}
    PV=d_TFV.
\end{equation*}

Il fattore di sconto $d_T$ ha un'interpretazione finanziaria
immediata: rappresenta il valore oggi di una unità monetaria che
verrà ricevuta con certezza al tempo $T$.
Ad esempio, se
$$    d_T=0.90,$$
una unità monetaria disponibile al tempo $T$ ha oggi un valore
pari a $0.90$.

Tasso di interesse e fattore di sconto rappresentano quindi due
modi differenti di esprimere la stessa relazione tra valori
monetari disponibili in momenti diversi.


\subsection{La struttura a termine dei tassi di interesse}

Finora abbiamo indicato il tasso di interesse con $r_T$,
specificando esplicitamente la scadenza $T$. Questa indicazione
è importante perché sui mercati finanziari non si utilizza
necessariamente lo stesso tasso per valutare pagamenti che si
verificano in date differenti.

In generale,
$$    r_1\neq r_2\neq\cdots\neq r_T.$$
Un pagamento che avverrà tra un anno può quindi essere scontato
utilizzando un tasso differente da quello utilizzato per un
pagamento che avverrà tra due, cinque o dieci anni.

La relazione che associa a ogni scadenza il corrispondente tasso
di interesse,

\begin{equation*}
    T\longmapsto r_T,
\end{equation*}

prende il nome di \emph{struttura a termine dei tassi di
interesse}.

La sua rappresentazione grafica è comunemente chiamata
\emph{curva dei rendimenti} o \emph{yield curve}.

I tassi $r_T$ utilizzati per scontare singoli pagamenti alle
diverse scadenze sono chiamati \emph{tassi spot} o
\emph{tassi zero-coupon}.


\subsection{Tassi spot e rendimento a scadenza}

È importante distinguere il tasso spot dal
\emph{rendimento a scadenza} di un'obbligazione
(\emph{yield to maturity}, YTM).

Il tasso spot $r_T$ è il tasso appropriato per scontare un
singolo pagamento che si verifica al tempo $T$:

\begin{equation*}
    PV=
    \frac{FV}{(1+r_T)^T}.
\end{equation*}

Consideriamo invece un'obbligazione che genera diversi flussi di
cassa
$$    CF_1,CF_2,\ldots,CF_T.$$
Il rendimento a scadenza $y$ è definito come l'unico tasso che,
applicato a tutti i flussi di cassa dell'obbligazione, rende il
loro valore attuale uguale al prezzo di mercato $P$:
$$
P=
    \sum_{t=1}^{T}
    \frac{CF_t}{(1+y)^t}.
    \label{eq:ytm}
$$
La differenza è quindi sostanziale.

Nel calcolo dello YTM utilizziamo lo stesso tasso $y$ per tutti
i flussi di cassa dell'obbligazione. Nella struttura a termine,
invece, ogni scadenza può essere associata a un proprio tasso
spot:
$$
    r_1,r_2,\ldots,r_T.
$$
Il rendimento a scadenza sintetizza quindi in un unico numero il
rendimento associato all'intera obbligazione, mentre i tassi spot
descrivono separatamente il valore temporale del denaro alle
diverse scadenze.

Poiché il nostro obiettivo è costruire la struttura a termine,
siamo interessati ai tassi spot $r_t$ e ai corrispondenti fattori
di sconto $d_t$, e non direttamente allo YTM dei singoli titoli.


\subsection{Valutazione di un'obbligazione}

Consideriamo ora un'obbligazione che genera i flussi di cassa
$$
    CF_1,CF_2,\ldots,CF_T.
$$
Poiché i pagamenti avvengono in date differenti, ciascun flusso
deve essere valutato utilizzando il fattore di sconto
corrispondente alla propria scadenza.

Il prezzo dell'obbligazione è quindi
$$
    P=
    CF_1d_1+
    CF_2d_2+
    \cdots+
    CF_Td_T,
$$
oppure, in forma compatta,
\begin{equation}
    P=
    \sum_{t=1}^{T}CF_td_t.
    \label{eq:bond-discount-factors}
\end{equation}

Questa relazione è fondamentale per il problema che vogliamo
studiare.

Se conosciamo i fattori di sconto, possiamo infatti utilizzarli
per calcolare il valore attuale dei flussi di cassa e quindi il
prezzo dell'obbligazione.


\subsection{Un semplice esempio}

Consideriamo un'obbligazione con scadenza a tre anni, valore
nominale pari a $100$ e cedola annuale pari a $5$.

I flussi di cassa sono
$$
    CF_1=5,
    \qquad
    CF_2=5,
    \qquad
    CF_3=105.
$$
Nei primi due anni l'investitore riceve soltanto la cedola. Alla
scadenza riceve invece l'ultima cedola, pari a $5$, insieme al
rimborso del valore nominale, pari a $100$.

Il prezzo dell'obbligazione è pertanto
\begin{equation*}
    P=5d_1+5d_2+105d_3.
    \label{eq:bond-three-years}
\end{equation*}
L'esempio mette in evidenza un aspetto importante: il prezzo di
un'obbligazione con cedole dipende generalmente da più fattori di
sconto, perché i suoi flussi di cassa si verificano in date
differenti.


\subsection{Il problema inverso: dai prezzi ai fattori di sconto}

L'Equazione~\eqref{eq:bond-discount-factors} può essere utilizzata
in due direzioni.

Se conosciamo i fattori di sconto
$$
    d_1,d_2,\ldots,d_T,
$$
e conosciamo i flussi di cassa dell'obbligazione, possiamo
calcolarne il prezzo.

Il problema che incontriamo nella realtà è però spesso quello
opposto.
I prezzi delle obbligazioni sono osservabili sui mercati
finanziari. Anche i flussi di cassa contrattualmente previsti
dalle obbligazioni sono noti. I fattori di sconto, invece, non
sono direttamente osservabili.

Possiamo allora chiederci se sia possibile utilizzare le
informazioni osservabili per ricavare quelle non osservabili.

Supponiamo di osservare $N$ obbligazioni e di considerare $N$
diverse scadenze. Per ciascuna obbligazione $i$ possiamo scrivere

\begin{equation*}
    P_i=
    \sum_{t=1}^{N}CF_{it}d_t,
    \qquad
    i=1,\ldots,N,
    \label{eq:bond-general-system}
\end{equation*}

dove $CF_{it}$ rappresenta il flusso di cassa generato
dall'obbligazione $i$ alla scadenza $t$.

Ogni obbligazione fornisce quindi un'equazione. Osservando più
obbligazioni otteniamo un sistema di equazioni nelle incognite
$$
    d_1,d_2,\ldots,d_N.
$$
Il problema finanziario può quindi essere formulato nel modo
seguente:
\begin{quote}
è possibile ricavare la struttura a termine dei tassi di interesse
utilizzando i prezzi delle obbligazioni osservati sul mercato?
\end{quote}


\subsection{Un esempio con tre obbligazioni}

Consideriamo tre obbligazioni, indicate con $A$, $B$ e $C$, con
i seguenti flussi di cassa e prezzi:

\begin{center}
\begin{tabular}{lrrrr}
\hline
Obbligazione & $t=1$ & $t=2$ & $t=3$ & Prezzo\\
\hline
A & 105 & 0   & 0   & 100\\
B & 4   & 104 & 0   & 99\\
C & 5   & 5   & 105 & 98\\
\hline
\end{tabular}
\end{center}

Per l'obbligazione $A$ abbiamo

\begin{equation*}
    105d_1=100.
    \label{eq:bond-A}
\end{equation*}

Per l'obbligazione $B$:

\begin{equation*}
    4d_1+104d_2=99.
    \label{eq:bond-B}
\end{equation*}

Infine, per l'obbligazione $C$:

\begin{equation*}
    5d_1+5d_2+105d_3=98.
    \label{eq:bond-C}
\end{equation*}

Le tre incognite sono quindi i fattori di sconto

\begin{equation*}
    d=
    \begin{pmatrix}
        d_1\\
        d_2\\
        d_3
    \end{pmatrix}.
\end{equation*}


\subsection{Dal problema finanziario al sistema lineare}

Le tre equazioni precedenti possono essere raccolte in forma
matriciale:

\begin{equation*}
\begin{pmatrix}
105 & 0   & 0\\
4   & 104 & 0\\
5   & 5   & 105
\end{pmatrix}
\begin{pmatrix}
d_1\\
d_2\\
d_3
\end{pmatrix}
=
\begin{pmatrix}
100\\
99\\
98
\end{pmatrix}.
\label{eq:bond-matrix-example}
\end{equation*}

Definiamo

\begin{equation*}
C=
\begin{pmatrix}
105 & 0   & 0\\
4   & 104 & 0\\
5   & 5   & 105
\end{pmatrix}
\end{equation*}

come matrice dei flussi di cassa,

\begin{equation*}
d=
\begin{pmatrix}
d_1\\
d_2\\
d_3
\end{pmatrix}
\end{equation*}

come vettore dei fattori di sconto e

\begin{equation*}
P=
\begin{pmatrix}
100\\
99\\
98
\end{pmatrix}
\end{equation*}

come vettore dei prezzi osservati.

Il problema può quindi essere scritto nella forma compatta

\begin{equation*}
	Cd=P.
    \label{eq:CdP}
\end{equation*}

Abbiamo così trasformato il problema finanziario in un sistema di
equazioni lineari.

La corrispondenza con la notazione generale utilizzata nelle
sezioni precedenti è immediata:

\begin{equation*}
\begin{array}{ccc}
Ax=b
&\longleftrightarrow&
Cd=P.
\end{array}
\end{equation*}

La matrice dei coefficienti $A$ è ora la matrice dei flussi di
cassa $C$; il vettore delle incognite $x$ è il vettore dei fattori
di sconto $d$; il vettore dei termini noti $b$ è costituito dai
prezzi osservati $P$.

Questo passaggio è particolarmente importante perché mostra che
una tecnica numerica generale può essere utilizzata per risolvere
problemi appartenenti ad ambiti applicativi molto differenti.

\iffalse
\subsection{Bootstrapping della struttura a termine}

Nell'esempio appena considerato la matrice dei flussi di cassa è
triangolare inferiore:

\begin{equation*}
C=
\begin{pmatrix}
105 & 0   & 0\\
4   & 104 & 0\\
5   & 5   & 105
\end{pmatrix}.
\end{equation*}

Questa particolare struttura permette di determinare i fattori di
sconto uno dopo l'altro.

Dalla prima equazione,

\begin{equation*}
    105d_1=100,
\end{equation*}

otteniamo immediatamente

\begin{equation*}
    d_1=\frac{100}{105}.
    \label{eq:d1-bootstrap}
\end{equation*}

Una volta noto $d_1$, possiamo utilizzare la seconda equazione:

\begin{equation*}
    4d_1+104d_2=99.
\end{equation*}

Isolando $d_2$ otteniamo

\begin{equation*}
    d_2=
    \frac{99-4d_1}{104}.
    \label{eq:d2-bootstrap}
\end{equation*}

A questo punto conosciamo $d_1$ e $d_2$ e possiamo utilizzare la
terza equazione:

\begin{equation*}
    5d_1+5d_2+105d_3=98,
\end{equation*}

da cui

\begin{equation*}
    d_3=
    \frac{98-5d_1-5d_2}{105}.
    \label{eq:d3-bootstrap}
\end{equation*}

Abbiamo quindi ricavato progressivamente

\begin{equation*}
    d_1
    \longrightarrow
    d_2
    \longrightarrow
    d_3.
\end{equation*}

Questa procedura sequenziale prende il nome di
\emph{bootstrapping}.

Nel nostro esempio il bootstrapping è particolarmente semplice
perché la matrice $C$ è triangolare. Dal punto di vista numerico,
stiamo quindi utilizzando la stessa logica della sostituzione in
avanti incontrata nello studio dei sistemi triangolari.

Più in generale, la struttura dei flussi di cassa potrebbe non
permettere una soluzione sequenziale così immediata. Il problema
rimane tuttavia rappresentabile come

\begin{equation*}
    Cd=P
\end{equation*}

e può essere affrontato utilizzando le tecniche numeriche per la
soluzione dei sistemi lineari introdotte nelle sezioni precedenti.
\fi

\subsection{Dai fattori di sconto ai tassi spot}

La soluzione del sistema ci fornisce i fattori di sconto
$$
    d_1,d_2,\ldots,d_T.
$$

Il nostro obiettivo finale è però costruire la struttura a termine
dei tassi di interesse. Dobbiamo quindi trasformare ciascun fattore
di sconto nel corrispondente tasso spot.

Ricordiamo la relazione

$$
    d_t=
    \frac{1}{(1+r_t)^t}.
$$

Elevando entrambi i membri alla potenza $-1/t$ otteniamo
$$
    1+r_t=d_t^{-1/t},
$$

e quindi
$$
    r_t=d_t^{-1/t}-1.
    \label{eq:discount-to-spot}
$$

Una volta calcolato un fattore di sconto per ciascuna scadenza,
possiamo pertanto ottenere i corrispondenti tassi spot:

\begin{equation*}
\begin{array}{ccccccc}
d_1 & d_2 & d_3 & \cdots & d_T\\
\downarrow & \downarrow & \downarrow && \downarrow\\
r_1 & r_2 & r_3 & \cdots & r_T.
\end{array}
\end{equation*}
Associando ciascun tasso $r_t$ alla propria scadenza $t$ otteniamo
la struttura a termine
$$
    t\longmapsto r_t,
$$
che può essere rappresentata graficamente attraverso la curva dei
rendimenti.

\paragraph{La struttura computazionale del problema}
Possiamo infine riassumere il procedimento seguito.

I dati inizialmente disponibili sono i prezzi delle obbligazioni
e i rispettivi flussi di cassa:

\begin{equation*}
    C,\qquad P.
\end{equation*}

Da queste informazioni costruiamo il sistema

\begin{equation*}
    Cd=P.
\end{equation*}

La soluzione del sistema fornisce i fattori di sconto:

\begin{equation*}
    C,P
    \longrightarrow
    d.
\end{equation*}

Successivamente trasformiamo i fattori di sconto nei
corrispondenti tassi spot:

\begin{equation*}
    d_t
    \longrightarrow
    r_t=d_t^{-1/t}-1.
\end{equation*}

L'intero procedimento può quindi essere rappresentato come

\begin{equation*}
\boxed{
\begin{array}{c}
\text{prezzi e flussi di cassa}\\
\downarrow\\
Cd=P\\
\downarrow\\
\text{fattori di sconto}\\
\downarrow\\
r_t=d_t^{-1/t}-1\\
\downarrow\\
\text{tassi spot}\\
\downarrow\\
\text{curva dei rendimenti}
\end{array}}
\end{equation*}

Il problema finanziario della costruzione della struttura a termine
si riconduce quindi, nella sua parte centrale, allo stesso problema
numerico

\begin{equation*}
    Ax=b
\end{equation*}

incontrato nelle precedenti applicazioni del corso.

\section{Richiami: titoli di Stato italiani}
\label{subsec:titoli-stato-italiani}

Finora abbiamo studiato il problema della costruzione della struttura
a termine utilizzando obbligazioni teoriche, caratterizzate da flussi
di cassa particolarmente semplici.

L'obiettivo è ora applicare la stessa procedura a dati osservati sui
mercati finanziari. Prima di procedere con l'applicazione numerica è
quindi utile richiamare brevemente le principali caratteristiche dei
titoli di Stato italiani.

I titoli di Stato sono strumenti finanziari emessi dal Ministero
dell'Economia e delle Finanze per finanziare il debito pubblico.
Dal punto di vista dell'investitore, l'acquisto di un titolo comporta
il pagamento di un prezzo al momento dell'acquisto e il diritto a
ricevere uno o più flussi di cassa futuri.

La struttura di questi flussi non è però uguale per tutti i titoli.
Alcuni strumenti non prevedono il pagamento di cedole, mentre altri
generano pagamenti periodici. Esistono inoltre titoli nei quali le
cedole o il capitale rimborsato dipendono dall'andamento di un
determinato indice.

Queste differenze sono importanti per la nostra applicazione.
Per costruire la matrice dei flussi di cassa $C$ dobbiamo infatti
conoscere, per ciascun titolo, non soltanto il prezzo osservato, ma
anche l'ammontare e la data di ciascun pagamento futuro.


\subsection{Buoni Ordinari del Tesoro}

I \emph{Buoni Ordinari del Tesoro} (BOT) sono titoli a breve termine
che non prevedono il pagamento di cedole.

Il rendimento dell'investitore deriva dalla differenza tra il prezzo
pagato per acquistare il titolo e il valore rimborsato alla scadenza.

In forma semplificata, la struttura dei flussi può essere
rappresentata come
$$
-P_0
\qquad\longrightarrow\qquad
F_T
$$
dove $P_0$ rappresenta il prezzo pagato oggi e $F_T$ il valore
rimborsato alla scadenza.

Dal punto di vista della costruzione della struttura a termine,
questa caratteristica rende i BOT particolarmente semplici.

Se un titolo genera un unico pagamento $F_T$ alla scadenza $T$,
il suo prezzo soddisfa infatti

\begin{equation*}
    P_0=F_Td_T.
\end{equation*}

Il fattore di sconto relativo a quella scadenza può quindi essere
ricavato direttamente:
$$
	    d_T=\frac{P_0}{F_T}.
$$
I titoli privi di cedole costituiscono pertanto il caso più semplice
per osservare la relazione tra prezzo, fattore di sconto e tasso
spot.

\subsection{Buoni del Tesoro Poliennali}

I \emph{Buoni del Tesoro Poliennali} (BTP) sono titoli caratterizzati
da una scadenza più lunga e dal pagamento periodico di cedole.

Indicando con $c_t$ la cedola pagata alla data $t$ e con $F$ il
valore nominale rimborsato alla scadenza, la sequenza dei flussi di
cassa può essere rappresentata schematicamente come
$$
-P_0
\longrightarrow
c_1
\longrightarrow
c_2
\longrightarrow
\cdots
\longrightarrow
c_T+F
$$
Il prezzo del titolo dipende quindi dal valore attuale di tutti i
pagamenti futuri:

\begin{equation*}
    P_0
    =
    c_1d_1+c_2d_2+\cdots+(c_T+F)d_T.
    \label{eq:btp-price}
\end{equation*}

Questa struttura è precisamente quella che abbiamo incontrato
nell'esempio teorico precedente.

A differenza di un BOT, quindi, un singolo BTP non permette
generalmente di determinare direttamente un unico fattore di sconto:
il suo prezzo dipende contemporaneamente dai fattori di sconto
associati a tutte le date nelle quali il titolo genera un pagamento.

È proprio questa caratteristica a rendere naturale la costruzione
del sistema

\begin{equation*}
    Cd=P.
\end{equation*}


\subsection{Certificati di Credito del Tesoro}

I \emph{Certificati di Credito del Tesoro} (CCT) appartengono alla
categoria dei titoli a tasso variabile.

A differenza di un titolo a cedola fissa, l'ammontare delle cedole
future non è necessariamente noto al momento dell'emissione, ma
dipende dal parametro al quale il titolo è indicizzato.

La struttura generale dei flussi rimane
$$
-P_0
\longrightarrow
c_1
\longrightarrow
c_2
\longrightarrow
\cdots
\longrightarrow
c_T+F,
$$
ma le cedole
$$
    c_1,c_2,\ldots,c_T
$$
non sono tutte fissate anticipatamente nello stesso modo di un
titolo a tasso fisso.

Questa differenza è rilevante per la costruzione della matrice dei
flussi di cassa: per utilizzare un titolo nella procedura descritta
precedentemente dobbiamo essere in grado di determinare i pagamenti
associati alle diverse date.


\subsection{Titoli indicizzati all'inflazione}

Un'ulteriore categoria è costituita dai titoli il cui rendimento è
collegato all'andamento dell'inflazione.

In questi strumenti l'indicizzazione ha l'obiettivo di collegare
almeno una parte dei pagamenti ricevuti dall'investitore
all'evoluzione del livello dei prezzi.

Dal punto di vista della nostra applicazione, la caratteristica
importante è che i flussi di cassa futuri non possono essere
trattati nello stesso modo dei flussi di un'obbligazione nominale
a cedola fissa.

Per costruire una struttura a termine nominale mediante il sistema
$$
    Cd=P,
$$
è infatti particolarmente conveniente lavorare con strumenti per
i quali i flussi di cassa nominali futuri siano determinabili.

\subsection{Perché la tipologia del titolo è importante}

Il breve richiamo alle diverse tipologie di titoli di Stato permette
di comprendere un punto importante per la successiva applicazione.

Il prezzo di mercato, da solo, non è sufficiente per costruire la
struttura a termine. Per ogni titolo dobbiamo conoscere l'intera
sequenza dei flussi di cassa:
$$
    CF_{i1},CF_{i2},\ldots,CF_{iT}.
$$
Questi valori costituiscono le righe della matrice
\begin{equation*}
C=
\begin{pmatrix}
CF_{11} & CF_{12} & \cdots & CF_{1T}\\
CF_{21} & CF_{22} & \cdots & CF_{2T}\\
\vdots  & \vdots  &        & \vdots\\
CF_{N1} & CF_{N2} & \cdots & CF_{NT}
\end{pmatrix}.
\end{equation*}

Il vettore

\begin{equation*}
P=
\begin{pmatrix}
P_1\\
P_2\\
\vdots\\
P_N
\end{pmatrix}
\end{equation*}

contiene invece i prezzi osservati sul mercato.

Il problema finanziario diventa ancora una volta

\begin{equation*}
    Cd=P.
\end{equation*}

Nell'esempio teorico precedente avevamo costruito artificialmente
titoli con date dei pagamenti perfettamente ordinate e una matrice
dei flussi di cassa triangolare.

Quando utilizziamo titoli realmente negoziati sul mercato, la
situazione è inevitabilmente più articolata. I titoli possono avere
date di scadenza differenti, cedole differenti e calendari dei
pagamenti che non coincidono perfettamente.

Prima di poter applicare gli strumenti numerici studiati in precedenza
sarà quindi necessario organizzare opportunamente queste
informazioni.
Questo passaggio ci porterà dall'esempio teorico alla costruzione
della struttura a termine utilizzando dati finanziari reali.

\section{Un esempio utilizzando i BOT}

\subsection{Descrizione e risultati}
Elaboriamo in questo paragrafo uno script che ha l'obiettivo di ricavare i \textbf{fattori di sconto dai prezzi di mercato dei BOT} (Buoni Ordinari del Tesoro), utilizzando dati disponibili online.

Inizialmente vengono caricate alcune librerie di R necessarie per la gestione dei dati, la lettura delle pagine web e la manipolazione delle stringhe di testo. Attraverso una procedura di \textit{web scraping}, lo script accede al sito \textit{Simple Tools for Investors} ed estrae una tabella contenente informazioni sui titoli di Stato italiani.

Dalla tabella originale vengono selezionate le variabili di interesse: codice ISIN, descrizione del titolo, data di scadenza, prezzo di riferimento, rendimento (\textit{yield}) e mercato di negoziazione. I dati vengono inoltre convertiti nei formati appropriati, in particolare le date e i valori numerici.

Successivamente, i titoli vengono classificati automaticamente in diverse categorie (BOT, BTP Italia, BTP indicizzati all'inflazione, BTP Valore, BTP Futura, BTP Più, BTP Green e BTP tradizionali), sulla base della loro descrizione.

L'analisi viene limitata ai \textbf{BOT e ai BTP tradizionali} (\textit{plain vanilla}), che vengono ordinati per data di scadenza e separati in due distinti dataset. In questa parte dello script vengono analizzati esclusivamente i BOT. Ad esempio, i BOT che erano disponibili il 22 ottobre 2026 sono riportati nella seguente tabella.

\begin{table}
\begin{tabular}{llrr}
\toprule
descrizione & scadenza & prezzo & giorni\_scadenza\\
\midrule
BOT 13/11/2026 ZC & 2026-11-13 & 99.85 & 22\\
BOT 30/11/2026 ZC & 2026-11-30 & 99.73 & 39\\
BOT 14/12/2026 ZC & 2026-12-14 & 99.67 & 53\\
BOT 31/12/2026 ZC & 2026-12-31 & 99.50 & 70\\
BOT 14/01/2027 ZC & 2027-01-14 & 99.40 & 84\\
\addlinespace
BOT 29/01/2027 ZC & 2027-01-29 & 99.25 & 99\\
BOT 12/02/2027 ZC & 2027-02-12 & 99.14 & 113\\
BOT 12/03/2027 ZC & 2027-03-12 & 98.93 & 141\\
BOT 31/03/2027 ZC & 2027-03-31 & 98.73 & 160\\
BOT 14/04/2027 ZC & 2027-04-14 & 98.61 & 174\\
\addlinespace
BOT 14/05/2027 ZC & 2027-05-14 & 98.35 & 204\\
BOT 14/06/2027 ZC & 2027-06-14 & 98.05 & 235\\
BOT 14/07/2027 ZC & 2027-07-14 & 97.78 & 265\\
BOT 13/08/2027 ZC & 2027-08-13 & 97.48 & 295\\
BOT 14/09/2027 ZC & 2027-09-14 & 97.21 & 327\\
\addlinespace
BOT 14/10/2027 ZC & 2027-10-14 & 96.93 & 357\\
\bottomrule
\end{tabular}
\caption{BOT in essere al 22 ottobre 2026}
\label{tab:allbot}
\end{table}

Per ciascun BOT con un prezzo disponibile, viene calcolato il tempo residuo alla scadenza, espresso in anni, attraverso la formula:
\[
T=\frac{\text{Data di scadenza}-\text{Data di valutazione}}{365}.
\]
La data di valutazione coincide con il giorno in cui viene eseguito lo script. Vengono esclusi i titoli già scaduti.

Poiché i BOT sono \textbf{titoli zero-coupon}, ossia non pagano cedole periodiche ma prevedono un unico rimborso alla scadenza, il loro prezzo può essere utilizzato direttamente per ricavare il fattore di sconto relativo alla corrispondente scadenza.

Assumendo un valore nominale pari a 100, il fattore di sconto è calcolato come:
\[
D(0,T)=\frac{P(0,T)}{100}.
\]
dove $P(0,T)$ rappresenta il prezzo del BOT alla data di valutazione e $D(0,T)$ il fattore di sconto relativo alla scadenza $T$.

Il fattore di sconto rappresenta il valore attuale di un euro che sarà ricevuto alla scadenza. Ad esempio, se un BOT ha un prezzo pari a 98, il corrispondente fattore di sconto è $0{,}98$: ciò significa che il valore attuale di un euro da ricevere alla scadenza è pari a 98 centesimi.

Infine, lo script costruisce un \textbf{grafico dei fattori di sconto in funzione del tempo residuo alla scadenza}. Ogni punto rappresenta un BOT, con la scadenza espressa in anni sull'asse orizzontale e il corrispondente fattore di sconto sull'asse verticale.

La figura \ref{fig:fds} permette di osservare la struttura per scadenza dei fattori di sconto ricavati dai prezzi di mercato dei BOT, che costituisce un primo passo per la costruzione della \textit{curva dei fattori di sconto} e, successivamente, della \textit{struttura a termine dei tassi di interesse}.

\begin{figure}
\centering
\includegraphics[scale=0.4]{fig_fasttori_di_sconto_bot.pdf}
\caption{I fattori di sconto dei BOT al 22 ottobre 2026}
\label{fig:fds}
\end{figure}



\paragraph{Nota metodologica.}
Il calcolo $P(0,T)/100$ assume che il prezzo utilizzato sia effettivamente il prezzo pagato per acquistare il titolo, espresso per 100 euro di valore nominale, e che non vi siano altri flussi di cassa. Inoltre, i fattori di sconto così ottenuti sono osservazioni relative alle singole scadenze disponibili, non ancora una curva continua. Per ottenere una curva completa sarà necessario ricorrere a tecniche di interpolazione o ad altre procedure di stima.




\subsection{Web scraping e manipolazione dei dati con R}
\label{sec:web_scraping_dplyr}

L'analisi quantitativa richiede spesso di acquisire informazioni da fonti esterne e di trasformarle in una struttura adatta alle successive elaborazioni. I dati pubblicati su Internet, per esempio, non sono sempre disponibili come file CSV o fogli di calcolo: possono essere presentati direttamente all'interno di pagine web sotto forma di tabelle HTML. In questi casi è possibile ricorrere al \emph{web scraping}, cioè all'estrazione automatizzata di informazioni dal contenuto di una pagina web. Dopo l'acquisizione, è generalmente necessario selezionare le variabili, correggere i formati, classificare le osservazioni ed eliminare quelle non pertinenti.

Consideriamo come esempio uno script R che legge una tabella di titoli di Stato italiani dal sito Simple Tools for Investors, distingue alcune categorie di strumenti e utilizza le quotazioni dei BOT per costruire una serie di fattori di sconto. L'esempio permette di illustrare quattro pacchetti: \texttt{rvest}, per leggere pagine HTML; \texttt{dplyr}, per manipolare tabelle; \texttt{readr}, per convertire dati testuali in numeri; e \texttt{stringr}, per riconoscere particolari sequenze di caratteri.

\subsubsection{Caricamento dei pacchetti e operatore pipe}

All'inizio dello script vengono caricati i pacchetti necessari mediante \texttt{library()}:
\begin{verbatim}
library(dplyr)
library(rvest)
library(readr)
library(stringr)
\end{verbatim}
Il comando \texttt{library()} rende disponibili nella sessione corrente le funzioni del pacchetto, che deve essere stato precedentemente installato. In molti passaggi viene impiegato l'operatore \emph{pipe} nativo di R, \texttt{|>}. Esso passa il risultato dell'espressione alla sua sinistra alla funzione successiva, normalmente come primo argomento. Per esempio:
\begin{verbatim}
dati |> filter(prezzo > 90) |> arrange(scadenza)
\end{verbatim}
L'istruzione seleziona le righe con prezzo superiore a 90 e successivamente le ordina per scadenza. La scrittura con pipe rende più evidente la successione delle trasformazioni e riduce la necessità di creare oggetti intermedi.

\subsubsection{Acquisizione di tabelle HTML con \texttt{rvest}}

Una pagina web è normalmente descritta da un documento HTML, nel quale i diversi elementi sono identificati da appositi tag. Le tabelle, per esempio, sono contenute in elementi \texttt{table}. Il pacchetto \texttt{rvest} consente di leggere la struttura della pagina e di estrarne gli elementi di interesse.

Nello script l'indirizzo della pagina è assegnato all'oggetto \texttt{url}; la funzione \texttt{read\_html()} ne acquisisce e interpreta il contenuto:
\begin{verbatim}
page <- read_html(url)
\end{verbatim}
L'oggetto \texttt{page} non è ancora una tabella R: rappresenta il documento HTML. Per cercare tutte le tabelle presenti nella pagina si utilizza \texttt{html\_elements("table")}. La funzione \texttt{html\_table()} prova poi a convertire ciascun elemento individuato in una tabella utilizzabile in R:
\begin{verbatim}
tables <- page |>
  html_elements("table") |>
  html_table(fill = TRUE)
\end{verbatim}
L'argomento \texttt{fill = TRUE} permette di gestire tabelle HTML nelle quali alcune righe hanno un numero diverso di celle, completando dove necessario la struttura. Il risultato è una lista di tabelle, non un'unica tabella. Con l'istruzione
\begin{verbatim}
btp_table <- tables[[1]]
\end{verbatim}
si estrae il primo elemento della lista. La notazione \texttt{[[1]]} serve infatti a recuperare il contenuto di un singolo elemento. Questa scelta presuppone che la tabella desiderata sia effettivamente la prima: se la struttura del sito cambia, potrebbe essere necessario modificare l'indice o identificare la tabella in altro modo.

Il web scraping dipende quindi dalla struttura della pagina e dalla sua disponibilità. Per un utilizzo ripetuto è opportuno controllare che la tabella individuata contenga le colonne attese e rispettare le condizioni d'uso del sito.

\subsubsection{Selezione e trasformazione delle colonne con \texttt{dplyr}}

Una tabella importata dal web può contenere molte variabili, con nomi poco comodi da utilizzare e formati non adeguati ai calcoli. La funzione \texttt{transmute()} del pacchetto \texttt{dplyr} consente di costruire una nuova tabella mantenendo soltanto le colonne specificate e, contemporaneamente, di rinominarle o trasformarle. Nello script vengono create le variabili \texttt{isin}, \texttt{descrizione}, \texttt{scadenza}, \texttt{prezzo}, \texttt{yield} e \texttt{mercato}. Un estratto del codice è:
\begin{verbatim}
btp <- btp_table |>
  transmute(
    isin = `Codice ISIN`,
    descrizione = Descrizione,
    scadenza = as.Date(Datascadenza),
    prezzo = parse_number(
      `Prezzo diriferimento`,
      locale = locale(decimal_mark = ",")
    )
  )
\end{verbatim}
I nomi delle colonne originali che contengono spazi possono essere racchiusi tra accenti gravi, come \texttt{`Codice ISIN`}. La funzione \texttt{as.Date()} converte la scadenza in un oggetto di tipo data, se il formato originario è riconosciuto. \texttt{parse\_number()}, appartenente a \texttt{readr}, estrae invece un numero da una stringa. L'argomento \texttt{locale(decimal\_mark = ",")} specifica che la virgola è il separatore decimale: questo dettaglio è importante per i dati espressi secondo le convenzioni italiane. Nello script la stessa conversione è applicata al rendimento, indicato dalla colonna \texttt{Yield}.

È utile distinguere \texttt{transmute()} da \texttt{mutate()}: la prima restituisce le colonne create o esplicitamente selezionate, mentre la seconda aggiunge o modifica variabili conservando anche le altre colonne della tabella.

\subsubsection{Classificazione dei titoli con \texttt{mutate()}, \texttt{case\_when()} e \texttt{stringr}}

Le descrizioni dei titoli contengono informazioni utili per identificarne la tipologia. Per automatizzare questa classificazione, lo script crea una nuova colonna \texttt{tipo} combinando \texttt{mutate()} e \texttt{case\_when()}. Quest'ultima funzione valuta una sequenza di condizioni e assegna a ogni osservazione il valore corrispondente alla prima condizione soddisfatta.

Le condizioni sono formulate mediante \texttt{str\_detect()}, funzione del pacchetto \texttt{stringr} che restituisce \texttt{TRUE} quando una stringa contiene un certo modello testuale. Per esempio:
\begin{verbatim}
btp <- btp |>
  mutate(
    tipo = case_when(
      str_detect(descrizione, "^BOT ") ~ "BOT",
      str_detect(descrizione, "^BTP Italia") ~ "BTP Italia",
      str_detect(descrizione, "^BTP ") ~ "BTP plain vanilla",
      TRUE ~ "Altro"
    )
  )
\end{verbatim}
Questo è un esempio semplificato: lo script completo riconosce anche BTP indicizzati all'inflazione, BTP Valore, Futura, Più e Green. Il simbolo \texttt{\^} all'interno del modello indica l'inizio della stringa: \texttt{"\^BOT "} riconosce dunque una descrizione che comincia con ``BOT'' seguito da uno spazio. Si tratta di un'espressione regolare, cioè di una sintassi per descrivere modelli di testo. La funzione \texttt{regex("GREEN", ignore\_case = TRUE)} permette invece di cercare la parola ``GREEN'' senza distinguere maiuscole e minuscole.

L'ordine delle condizioni in \texttt{case\_when()} è importante. Una categoria generica, come i BTP ordinari, deve essere verificata dopo le categorie più specifiche, altrimenti potrebbe assorbire osservazioni destinate a queste ultime. La condizione finale \texttt{TRUE \textasciitilde{} "Altro"} attribuisce un'etichetta alle osservazioni che non soddisfano nessuna regola precedente.

Per verificare quali categorie sono presenti, lo script utilizza:
\begin{verbatim}
tipi_presenti <- levels(as.factor(btp$tipo))
print(tipi_presenti)
\end{verbatim}
\texttt{as.factor()} trasforma i valori della variabile in categorie, mentre \texttt{levels()} ne restituisce i livelli; \texttt{print()} li visualizza nella console.

\subsubsection{Filtraggio, ordinamento e separazione dei dati}

Una volta classificati i titoli, si possono selezionare soltanto quelli necessari all'analisi. \texttt{filter()} mantiene le righe che soddisfano una condizione logica; \texttt{arrange()} ordina le osservazioni; \texttt{select()} sceglie le colonne da visualizzare o conservare. Nello script:
\begin{verbatim}
titoli_selezionati <- btp |>
  filter(tipo %in% c("BOT", "BTP plain vanilla")) |>
  arrange(scadenza)
\end{verbatim}
L'operatore \texttt{\%in\%} controlla se ciascun valore di \texttt{tipo} appartiene all'insieme indicato con \texttt{c()}. I titoli selezionati sono poi ordinati per scadenza crescente. Una successiva \texttt{select()} mostra le variabili più rilevanti, senza modificare l'oggetto \texttt{titoli\_selezionati} se il risultato non viene assegnato.

Infine, i due gruppi vengono separati in altrettante tabelle:
\begin{verbatim}
bot <- titoli_selezionati |>
  filter(tipo == "BOT")

btp_plain <- titoli_selezionati |>
  filter(tipo == "BTP plain vanilla")
\end{verbatim}
L'operatore \texttt{==} verifica l'uguaglianza, mentre \texttt{=} è utilizzato, tra l'altro, per assegnare valori agli argomenti delle funzioni.

\subsubsection{Gestione dei valori mancanti e calcolo delle scadenze}

Prima di procedere ai calcoli finanziari, lo script elimina i BOT per i quali manca il prezzo:
\begin{verbatim}
bot <- bot |> filter(!is.na(prezzo))
\end{verbatim}
In R i valori mancanti sono rappresentati da \texttt{NA}. La funzione \texttt{is.na()} li individua e l'operatore \texttt{!} nega la condizione; pertanto vengono conservate soltanto le osservazioni con un prezzo disponibile.

La data di valutazione viene fissata al giorno in cui lo script è eseguito, mediante \texttt{Sys.Date()}. Per ogni BOT si calcola quindi il numero di giorni residui e lo si converte in anni:
\begin{verbatim}
data_valutazione <- Sys.Date()
bot <- bot |>
  mutate(
    giorni_scadenza = as.numeric(scadenza - data_valutazione),
    T = giorni_scadenza / 365
  ) |>
  filter(T > 0)
\end{verbatim}
La sottrazione tra due date produce una durata; \texttt{as.numeric()} la esprime numericamente in giorni. Dividendo per 365 si ottiene una misura approssimata della vita residua in anni, secondo la convenzione adottata nello script. L'ultimo filtro esclude i titoli già scaduti o con scadenza coincidente con la data di valutazione. Poiché \texttt{Sys.Date()} cambia da un giorno all'altro, anche le scadenze residue cambiano quando lo script viene rieseguito.

\subsubsection{Dai prezzi dei BOT ai fattori di sconto}

Un fattore di sconto relativo alla scadenza $T$, indicato con $D(T)$, rappresenta il valore attuale di un'unità monetaria pagata alla data $T$. Per uno strumento a sconto che non distribuisce cedole e rimborsa il valore nominale $N$ alla scadenza, la relazione tra prezzo e fattore di sconto è
$$
P(T)=N D(T).
$$
Di conseguenza,
$$
D(T)=\frac{P(T)}{N}.
$$
Nel caso di un prezzo espresso in percentuale del valore nominale, ponendo $N=100$, si ottiene
$$
D(T)=\frac{P(T)}{100}.
$$
Lo script applica questa relazione mediante una nuova variabile:
\begin{verbatim}
bot <- bot |>
  mutate(discount_factor = prezzo / 100)
\end{verbatim}
Questa operazione interpreta il prezzo riportato come prezzo di uno zero-coupon con rimborso pari a 100. È una semplificazione utile per l'esercizio didattico: in un'applicazione finanziaria occorre verificare la definizione esatta del prezzo acquisito, le convenzioni di mercato e gli eventuali aggiustamenti necessari.

\subsubsection{Rappresentazione grafica ed esportazione}

Dopo aver costruito i fattori di sconto, lo script li rappresenta in funzione della vita residua utilizzando la funzione \texttt{plot()} di R base:
\begin{verbatim}
plot(
  bot$T,
  bot$discount_factor,
  xlab = "Maturity (anni)",
  ylab = "Discount factor",
  main = "Discount factors ricavati dai BOT",
  pch = 19
)
\end{verbatim}
Il simbolo \texttt{\$} permette di accedere a una colonna di una tabella: \texttt{bot\$T} contiene le scadenze e \texttt{bot\$discount\_factor} i corrispondenti fattori di sconto. Gli argomenti \texttt{xlab}, \texttt{ylab} e \texttt{main} specificano rispettivamente le etichette degli assi e il titolo del grafico; \texttt{pch = 19} sceglie un simbolo pieno per i punti.

Lo script contiene inoltre, come istruzioni commentate, i comandi \texttt{pdf()}, \texttt{dev.off()} e \texttt{write.csv()}. I primi due permettono di aprire e chiudere un dispositivo grafico PDF, salvando il grafico in un file; \texttt{write.csv()} esporta una tabella in formato CSV. L'opzione \texttt{row.names = FALSE} evita di scrivere nel file una colonna aggiuntiva con i numeri delle righe.

\subsubsection{Considerazioni conclusive}

L'esempio mostra una tipica sequenza di lavoro con dati acquisiti dal web: lettura della pagina HTML, individuazione della tabella, conversione dei tipi di dato, classificazione delle osservazioni, selezione dei titoli rilevanti e costruzione di nuove variabili per l'analisi economica. Il vantaggio di questa impostazione è la riproducibilità: le operazioni vengono descritte in un unico script e possono essere ripetute quando i dati disponibili sulla pagina vengono aggiornati.

La riproducibilità, tuttavia, richiede alcuni controlli. La struttura HTML può cambiare, i nomi delle colonne possono essere modificati e le date o i numeri possono essere pubblicati in formati differenti. Inoltre, la correttezza delle trasformazioni statistiche non garantisce da sola la correttezza dell'interpretazione economica dei dati. È quindi opportuno verificare sia il risultato dell'importazione sia il significato finanziario delle variabili prima di utilizzare i fattori di sconto per ulteriori elaborazioni, come l'interpolazione della curva per scadenze non osservate.

\section{Interpolazione numerica}
\label{sec:interpolazione}

Nella sezione precedente abbiamo visto che i prezzi dei titoli di Stato
possono essere utilizzati per ricavare informazioni sulla struttura a
termine dei fattori di sconto. In particolare, i BOT sono strumenti
particolarmente semplici perché generano un unico pagamento alla
scadenza. Se il valore nominale è pari a 100, dal prezzo di un BOT
possiamo ricavare direttamente il fattore di sconto:

\begin{equation*}
    d_T = \frac{P_T}{100}.
\end{equation*}

Questo risultato sembra fornire una soluzione immediata al nostro
problema. Esiste tuttavia una difficoltà importante quando passiamo
dall'esempio teorico ai dati realmente osservati sul mercato.

I BOT disponibili non hanno una scadenza per ogni possibile data
futura. Essi permettono quindi di conoscere direttamente la funzione
di sconto soltanto in corrispondenza di un insieme discreto di
scadenze.

Supponiamo, ad esempio, di conoscere i fattori di sconto

\begin{equation*}
    d(T_1), d(T_2), \ldots, d(T_n)
\end{equation*}

corrispondenti alle scadenze

\begin{equation*}
    T_1,T_2,\ldots,T_n.
\end{equation*}

Potremmo però avere la necessità di valutare un pagamento che avviene
a una data $T^*$ che non coincide con nessuna di queste scadenze.

Il problema diventa allora

\begin{equation*}
\boxed{
\begin{array}{c}
\text{conosciamo } d(T_1),d(T_2),\ldots,d(T_n)\\[2mm]
\downarrow\\[2mm]
\text{vogliamo determinare } d(T^*)
\end{array}}
\end{equation*}

per una scadenza $T^*$ non direttamente osservata.

Questo è un problema di \emph{interpolazione}.

L'interpolazione consiste nel costruire una funzione che passa
attraverso un insieme di punti noti e nell'utilizzare tale funzione
per stimare il valore corrispondente a punti intermedi.

Nel nostro caso i punti osservati sono costituiti dalle coppie

\begin{equation*}
    (T_i,d_i),
\end{equation*}

dove $T_i$ rappresenta la maturity e $d_i$ il corrispondente fattore
di sconto.


\subsection{Interpolazione lineare}

Consideriamo inizialmente il caso più semplice. Supponiamo di conoscere
due punti della funzione di sconto:

\begin{equation*}
    (T_1,d_1)
    \qquad\text{e}\qquad
    (T_2,d_2),
\end{equation*}

e di voler stimare il fattore di sconto relativo a una maturity
intermedia $T^*$, con

\begin{equation*}
    T_1<T^*<T_2.
\end{equation*}

L'ipotesi più semplice consiste nell'assumere che, nell'intervallo
compreso tra $T_1$ e $T_2$, la funzione di sconto possa essere
approssimata mediante una retta:

\begin{equation*}
    d(T)=a+bT.
    \label{eq:linear-interpolation}
\end{equation*}

I coefficienti $a$ e $b$ non sono noti. Tuttavia conosciamo due punti
attraverso i quali la retta deve necessariamente passare.

Imponendo

\begin{equation*}
    d(T_1)=d_1
\end{equation*}

otteniamo

\begin{equation*}
    a+bT_1=d_1.
\end{equation*}

Analogamente, imponendo

\begin{equation*}
    d(T_2)=d_2
\end{equation*}

otteniamo

\begin{equation*}
    a+bT_2=d_2.
\end{equation*}

Abbiamo quindi due equazioni nelle due incognite $a$ e $b$:

\begin{equation*}
\begin{cases}
a+bT_1=d_1,\\
a+bT_2=d_2.
\end{cases}
\end{equation*}

Anche il problema dell'interpolazione può quindi essere ricondotto
alla soluzione di un sistema lineare.


\subsubsection{Interpolazione lineare come sistema di equazioni}

Le due equazioni precedenti possono essere scritte in forma
matriciale:

\begin{equation*}
\begin{pmatrix}
1 & T_1\\
1 & T_2
\end{pmatrix}
\begin{pmatrix}
a\\
b
\end{pmatrix}
=
\begin{pmatrix}
d_1\\
d_2
\end{pmatrix}.
\label{eq:linear-interpolation-system}
\end{equation*}

Ritroviamo ancora una volta la struttura generale

\begin{equation*}
    Ax=b.
\end{equation*}

In questo caso

\begin{equation*}
A=
\begin{pmatrix}
1&T_1\\
1&T_2
\end{pmatrix},
\qquad
x=
\begin{pmatrix}
a\\
b
\end{pmatrix},
\qquad
b=
\begin{pmatrix}
d_1\\
d_2
\end{pmatrix}.
\end{equation*}

Risolvendo il sistema determiniamo i coefficienti $a$ e $b$ della
retta. A quel punto possiamo calcolare il fattore di sconto
corrispondente a qualsiasi maturity $T^*$ compresa nell'intervallo:

\begin{equation*}
    d(T^*)=a+bT^*.
\end{equation*}

È importante osservare il ruolo svolto dal sistema lineare.
Le incognite non sono direttamente i fattori di sconto che vogliamo
interpolare, ma i coefficienti della funzione utilizzata per
rappresentare la relazione tra maturity e fattore di sconto.

Una volta determinata questa funzione, essa può essere valutata in
qualunque punto dell'intervallo considerato.


\subsubsection{Un esempio con dati finanziari reali}

Il problema appena descritto emerge direttamente nei dati relativi
ai titoli di Stato.

Consideriamo il titolo

\begin{center}
\texttt{BTP 01/12/2026 1,25\%}.
\end{center}

Per valutare il pagamento che avviene alla scadenza del titolo abbiamo
bisogno del fattore di sconto relativo al 1 dicembre 2026.

Supponiamo però che nel dataset non sia presente un BOT con questa
esatta scadenza. I due BOT con scadenze più vicine sono:

\begin{center}
\begin{tabular}{lll}
\hline
Titolo & Scadenza & Fattore di sconto\\
\hline
BOT 13/11/2026 ZC & 13/11/2026 & $d_2$\\
BOT 14/12/2026 ZC & 14/12/2026 & $d_1$\\
\hline
\end{tabular}
\end{center}

La data del pagamento del BTP si trova tra le due date:

\begin{equation*}
13/11/2026
<
01/12/2026
<
14/12/2026.
\end{equation*}

Possiamo quindi utilizzare i fattori di sconto dei BOT per
costruire una retta di interpolazione e valutare tale retta in
corrispondenza del 1 novembre 2026.

In termini generali, il procedimento è

\begin{equation*}
\boxed{
(T_1,d_1),(T_2,d_2)
\longrightarrow
\begin{pmatrix}
1&T_1\\
1&T_2
\end{pmatrix}
\begin{pmatrix}
a\\
b
\end{pmatrix}
=
\begin{pmatrix}
d_1\\
d_2
\end{pmatrix}
\longrightarrow
d(T^*)
}.
\end{equation*}

L'esempio mostra perché l'interpolazione è necessaria nella pratica:
le date dei pagamenti dei BTP non coincidono necessariamente con le
scadenze dei BOT dai quali possiamo osservare direttamente i fattori
di sconto.


\subsubsection{Dall'interpolazione lineare all'interpolazione polinomiale}

L'interpolazione lineare utilizza due punti osservati e costruisce una
retta che li collega:

\begin{equation*}
    d(T)=a_0+a_1T.
\end{equation*}

Possiamo però generalizzare questa idea.

Supponiamo, ad esempio, di conoscere tre punti:

\begin{equation*}
    (T_1,d_1),\qquad
    (T_2,d_2),\qquad
    (T_3,d_3).
\end{equation*}

Invece di utilizzare una retta possiamo cercare un polinomio di
secondo grado:

\begin{equation*}
    d(T)=a_0+a_1T+a_2T^2.
\end{equation*}

Imponendo che il polinomio passi attraverso i tre punti otteniamo

\begin{equation*}
\begin{cases}
a_0+a_1T_1+a_2T_1^2=d_1,\\
a_0+a_1T_2+a_2T_2^2=d_2,\\
a_0+a_1T_3+a_2T_3^2=d_3.
\end{cases}
\end{equation*}

Abbiamo quindi tre equazioni nelle tre incognite

\begin{equation*}
    a_0,\qquad a_1,\qquad a_2.
\end{equation*}

La stessa idea può essere estesa a un numero arbitrario di punti.


\subsection{Interpolazione polinomiale}

Consideriamo $n$ BOT dai quali abbiamo ricavato $n$ punti della
funzione di sconto:

\begin{equation*}
    (T_1,d_1),(T_2,d_2),\ldots,(T_n,d_n).
\end{equation*}

Cerchiamo ora un polinomio di grado $n-1$:

\begin{equation*}
    d(T)
    =
    a_0+a_1T+a_2T^2+\cdots+a_{n-1}T^{n-1}.
    \label{eq:polynomial-interpolation}
\end{equation*}

Il polinomio contiene esattamente $n$ coefficienti incogniti:

\begin{equation*}
    a_0,a_1,\ldots,a_{n-1}.
\end{equation*}

Per determinarli imponiamo che il polinomio passi attraverso tutti i
punti osservati:

\begin{equation*}
    d(T_i)=d_i,
    \qquad i=1,\ldots,n.
\end{equation*}

Otteniamo quindi il sistema

\begin{equation*}
\begin{cases}
a_0+a_1T_1+a_2T_1^2+\cdots+a_{n-1}T_1^{n-1}=d_1,\\
a_0+a_1T_2+a_2T_2^2+\cdots+a_{n-1}T_2^{n-1}=d_2,\\
\hspace{2cm}\vdots\\
a_0+a_1T_n+a_2T_n^2+\cdots+a_{n-1}T_n^{n-1}=d_n.
\end{cases}
\end{equation*}

Ancora una volta, un problema apparentemente differente può essere
ricondotto alla soluzione di un sistema di equazioni lineari.


\subsubsection{La matrice di Vandermonde}

Il sistema dell'interpolazione polinomiale può essere scritto in
forma matriciale:

\begin{equation*}
\underbrace{
\begin{pmatrix}
1&T_1&T_1^2&\cdots&T_1^{n-1}\\
1&T_2&T_2^2&\cdots&T_2^{n-1}\\
\vdots&\vdots&\vdots&&\vdots\\
1&T_n&T_n^2&\cdots&T_n^{n-1}
\end{pmatrix}
}_{A}
\underbrace{
\begin{pmatrix}
a_0\\
a_1\\
\vdots\\
a_{n-1}
\end{pmatrix}
}_{x}
=
\underbrace{
\begin{pmatrix}
d_1\\
d_2\\
\vdots\\
d_n
\end{pmatrix}
}_{b}.
\label{eq:vandermonde-system}
\end{equation*}

La matrice

\begin{equation*}
A=
\begin{pmatrix}
1&T_1&T_1^2&\cdots&T_1^{n-1}\\
1&T_2&T_2^2&\cdots&T_2^{n-1}\\
\vdots&\vdots&\vdots&&\vdots\\
1&T_n&T_n^2&\cdots&T_n^{n-1}
\end{pmatrix}
\end{equation*}

prende il nome di \emph{matrice di Vandermonde}.

Il problema dell'interpolazione assume quindi nuovamente la forma

\begin{equation*}
    Ax=b.
\end{equation*}

Questa formulazione mostra ancora una volta la generalità degli
strumenti numerici studiati nel corso: la stessa struttura matematica
compare nella soluzione di problemi economici, ambientali e
finanziari e, ora, anche nella costruzione di una funzione interpolante.

\subsubsection{Interpolazione con dati sui BOT}

Effettuiamo l'interpolazione polinomiale con sei dei BOT elencati nella tabella \ref{tab:allbot}:

\begin{verbatim}
1 BOT 13/11/2026 ZC
2 BOT 31/12/2026 ZC
3 BOT 12/02/2027 ZC
4 BOT 14/04/2027 ZC
5 BOT 14/07/2027 ZC
6 BOT 14/09/2027 ZC
\end{verbatim}

Il polinomio interpolante vien mostrato nella figura \ref{fig:interp_bot_6}. La figura mostra che possiamo utilizzare i parametri stimati per ottenere il fattore di sconto per qualsiasi scadenza che cada nell'intervallo considerato.

\begin{figure}
\centering
\includegraphics[scale=0.4]{fig_interp_bot6.pdf}
\caption{interpolazione polinomiale}
\label{fig:interp_bot_6}
\end{figure}


\subsubsection{Perché non utilizzare semplicemente molti punti?}

A prima vista sembrerebbe naturale utilizzare il maggior numero
possibile di osservazioni.

Se disponiamo di $n$ BOT possiamo costruire un polinomio di grado
$n-1$ che passa esattamente attraverso tutti gli $n$ punti osservati.

Tuttavia, aumentando il numero di punti aumenta anche il grado del
polinomio e la matrice di Vandermonde contiene potenze sempre più
elevate delle maturities:

\begin{equation*}
    1,\quad T,\quad T^2,\quad T^3,\quad\ldots,\quad T^{n-1}.
\end{equation*}

Questo può generare un problema numerico importante: la matrice può
diventare \emph{mal condizionata}.

Ciò significa che piccole variazioni nei dati possono produrre
variazioni molto più grandi nella soluzione del sistema.

Questo aspetto è particolarmente rilevante quando lavoriamo con dati
reali. Prezzi, rendimenti e fattori di sconto sono infatti quantità
misurate numericamente e possono essere soggetti ad arrotondamenti,
errori di misura o piccole variazioni di mercato.


\subsubsection{Il numero di condizionamento}

Consideriamo nuovamente un sistema lineare

\begin{equation*}
    Ax=b.
\end{equation*}

Una misura della sensibilità numerica del problema è fornita dal
\emph{numero di condizionamento} della matrice $A$, definito come

\begin{equation*}
    \kappa(A)=\|A\|\,\|A^{-1}\|
    \label{eq:condition-number}
\end{equation*}

dove $\|A\|$ è una norma matriciale (le più utilizzate sono la norma 1, 2 o $\infty$). La norma 1 ($\|A\|_1$) di una matrice è il massimo della somma dei valori assoluti per colonna. La norma $\infty$ ($\|A\|_\infty$) di una matrice è il massimo della somma dei valori assoluti per riga. La norma 2 è più complessa da calcolare, tuttavia, possiamo dare un limite superiore a questa norma utilizzando la seguente disuguaglianza
\[
\|A\|_2 \leq \sqrt{\|A\|_1 \|A\|_\infty}
\] 
come vediamo, il calcolo del numero di condizionamento richiede il calcolo dell'inversa. Per matrici grandi, esistono metodi di stima del numero di condizionamento che evitano il calcolo dell'inversa. 

Il numero di condizionamento ci dice quanto la soluzione del sistema
può essere sensibile a piccole perturbazioni nei dati.

Se il vettore dei termini noti subisce una perturbazione

\begin{equation*}
    b\longrightarrow b+\Delta b,
\end{equation*}

anche la soluzione cambia:

\begin{equation*}
    x\longrightarrow x+\Delta x.
\end{equation*}

Una relazione importante è

\begin{equation*}
    \frac{\|\Delta x\|}{\|x\|}
    \leq
    \kappa(A)
    \frac{\|\Delta b\|}{\|b\|}.
    \label{eq:condition-bound}
\end{equation*}

Questa espressione permette di comprendere il significato economico
e numerico del numero di condizionamento.

In questo caso, tutttavia, le norme che compaiono sono norme vettoriali. La norma 1 è la somma dei valori assoluti del vettore. La norma $\infty$ è il massimo dei componenti del vettore in volore assoluto. La norma 2 è la norma euclidea. Nel caso dei vettori abbiamo che il calore della norma 2 è compreso tra il valore della norma 1 e della norma $\infty$.

Se $\kappa(A)$ è piccolo, una piccola variazione nei dati produce
generalmente una piccola variazione nella soluzione. Se invece
$\kappa(A)$ è molto grande, una perturbazione apparentemente
trascurabile nei dati può essere fortemente amplificata.


\subsubsection{Interpretazione del condizionamento}

In termini qualitativi possiamo distinguere tre situazioni:

\begin{equation*}
\begin{array}{lll}
\kappa(A)\approx 1
&\Rightarrow&
\text{sistema ben condizionato},\\[2mm]
\kappa(A)\gg 1
&\Rightarrow&
\text{sistema mal condizionato},\\[2mm]
\kappa(A)=\infty
&\Rightarrow&
\text{matrice singolare}.
\end{array}
\end{equation*}

È importante distinguere il problema dell'esistenza della soluzione
dal problema della sua stabilità numerica.

Una matrice può infatti essere invertibile e il sistema può avere una
soluzione unica, ma la soluzione può essere estremamente sensibile a
piccole variazioni nei dati.

Pertanto,

\begin{quote}
avere una soluzione matematica unica non significa necessariamente
avere un problema numericamente ben posto.
\end{quote}

Questa osservazione è particolarmente importante per l'interpolazione
polinomiale. Utilizzare un numero crescente di punti permette di
costruire un polinomio che passa esattamente attraverso tutte le
osservazioni, ma può contemporaneamente peggiorare le proprietà
numeriche del problema.

\subsubsection{Il numero di condizionamento per l'esempio dei BOT}

Calcoliamo il numero di condizionamento utilizzando le dunzioni messe a disposizione dal nostro software.
In particolare confrontiamo una stima fatta con 6 titoli, con una in cui consideriamo 15 titoli. Il risultato è il seguente

\begin{verbatim}
       6 BOT       15 BOT 
1.598173e+04 5.564928e+13
\end{verbatim}

Il numero di condizionamento è sensibilmente più alto quando consideriamo 15 titoli.

Nasce quindi una nuova domanda: è possibile utilizzare molti punti
della funzione di sconto senza ricorrere a un unico polinomio di
grado elevato?

La risposta conduce a una diversa tecnica di interpolazione, basata
sull'utilizzo di polinomi di basso grado definiti localmente su
intervalli successivi: le \emph{spline}.

\subsection{Interpolazione mediante spline cubiche}
\label{sec:spline}

Nella sezione precedente abbiamo visto che l'interpolazione polinomiale
permette di costruire una funzione che passa esattamente attraverso un
insieme di punti osservati. Se disponiamo di $n$ punti, possiamo infatti
costruire un polinomio di grado $n-1$ determinandone i coefficienti
attraverso la soluzione di un sistema lineare basato sulla matrice di
Vandermonde.

Questa procedura presenta tuttavia un limite importante. All'aumentare
del numero di punti aumenta anche il grado del polinomio e la matrice
di Vandermonde può diventare mal condizionata. Di conseguenza, piccole
perturbazioni nei dati possono produrre variazioni rilevanti nei
coefficienti del polinomio e quindi nella funzione interpolata.

Sorge allora una domanda naturale: possiamo utilizzare molti punti
osservati senza costruire un unico polinomio di grado elevato?

Una possibile risposta consiste nell'abbandonare l'idea di utilizzare
un unico polinomio sull'intero intervallo e costruire invece diversi
polinomi di basso grado, ciascuno valido soltanto localmente.

Questa è l'idea fondamentale alla base delle \emph{spline}.


\subsubsection{L'idea delle spline}

Supponiamo di disporre di $n$ punti osservati

\begin{equation*}
    (T_1,d_1),(T_2,d_2),\ldots,(T_n,d_n),
\end{equation*}

ordinati in modo che

\begin{equation*}
    T_1<T_2<\cdots<T_n.
\end{equation*}

Le maturities $T_1,\ldots,T_n$ vengono chiamate \emph{nodi}.

Nell'interpolazione polinomiale globale cercheremmo un unico
polinomio

\begin{equation*}
    p_{n-1}(T)
\end{equation*}

di grado $n-1$ che passi attraverso tutti gli $n$ punti.

Con una spline adottiamo invece una strategia differente. Dividiamo
l'intervallo complessivo nei sottointervalli

\begin{equation*}
[T_1,T_2],\quad
[T_2,T_3],\quad
\ldots,\quad
[T_{n-1},T_n]
\end{equation*}

e associamo a ciascun intervallo una funzione polinomiale distinta.

In altre parole, non utilizziamo più un'unica funzione globale, ma una
funzione definita a tratti:

\begin{equation*}
S(T)=
\begin{cases}
S_1(T), & T_1\leq T\leq T_2,\\
S_2(T), & T_2\leq T\leq T_3,\\
\vdots & \\
S_{n-1}(T), & T_{n-1}\leq T\leq T_n.
\end{cases}
\end{equation*}

Il problema consiste quindi nel determinare i diversi polinomi
$S_i(T)$ e nel collegarli in modo tale che, nel loro insieme,
costituiscano una curva sufficientemente regolare.


\subsubsection{La spline cubica}

Una scelta molto comune consiste nell'utilizzare polinomi di terzo
grado. In questo caso si parla di \emph{spline cubica}.

Per ciascun intervallo $[T_i,T_{i+1}]$ definiamo

\begin{equation*}
    S_i(T)
    =
    a_i+b_iT+c_iT^2+e_iT^3,
    \qquad i=1,\ldots,n-1.
    \label{eq:cubic-spline}
\end{equation*}

Ogni intervallo possiede quindi quattro coefficienti:

\begin{equation*}
    a_i,\quad b_i,\quad c_i,\quad e_i.
\end{equation*}

Poiché $n$ nodi individuano $n-1$ intervalli, il numero complessivo
dei coefficienti da determinare è

\begin{equation*}
    4(n-1).
    \label{eq:spline-number-coefficients}
\end{equation*}

Nel nostro esempio finanziario disponiamo di 15 BOT. Abbiamo quindi

\begin{equation*}
    n=15
\end{equation*}

nodi e

\begin{equation*}
    n-1=14
\end{equation*}

intervalli.

Il numero totale dei coefficienti incogniti è pertanto

\begin{equation*}
    4(15-1)=56.
\end{equation*}

La costruzione diretta della spline richiede dunque la determinazione
di 56 coefficienti.

Il problema può sembrare notevolmente più complesso
dell'interpolazione lineare. Vedremo però che le proprietà che
desideriamo imporre alla spline generano precisamente le equazioni
necessarie per determinare questi coefficienti.


\subsubsection{Prima condizione: interpolazione dei nodi}

La prima proprietà è la stessa che caratterizza qualsiasi funzione
interpolante: la spline deve passare attraverso i punti osservati.

Consideriamo il polinomio $S_i(T)$ definito nell'intervallo

\begin{equation*}
    [T_i,T_{i+1}].
\end{equation*}

Esso deve assumere il valore $d_i$ nel primo estremo e il valore
$d_{i+1}$ nel secondo estremo. Dobbiamo quindi imporre

\begin{equation*}
    S_i(T_i)=d_i
\end{equation*}

e

\begin{equation*}
    S_i(T_{i+1})=d_{i+1},
\end{equation*}

per

\begin{equation*}
    i=1,\ldots,n-1.
\end{equation*}

Ogni intervallo fornisce quindi due equazioni. Poiché gli intervalli
sono $n-1$, otteniamo complessivamente

\begin{equation*}
    2(n-1)
\end{equation*}

equazioni.

Nel caso dei 15 BOT:

\begin{equation*}
    2(15-1)=28.
\end{equation*}

Le condizioni di interpolazione forniscono dunque 28 delle 56
equazioni necessarie.


\subsubsection{Seconda condizione: continuità della derivata prima}

Il semplice passaggio attraverso i nodi non è sufficiente.

Potremmo infatti costruire polinomi che si incontrano nei nodi ma che
presentano bruschi cambiamenti di pendenza. La funzione risultante
sarebbe continua, ma mostrerebbe degli ``spigoli'' nei punti di
raccordo.

Per ottenere una curva regolare richiediamo quindi che, in ciascun
nodo interno, i due polinomi adiacenti abbiano la stessa derivata
prima.

Nel nodo $T_i$ imponiamo pertanto

\begin{equation*}
    S'_{i-1}(T_i)=S'_i(T_i),
    \qquad i=2,\ldots,n-1.
    \label{eq:spline-first-derivative}
\end{equation*}

Partendo da

\begin{equation*}
    S_i(T)=a_i+b_iT+c_iT^2+e_iT^3,
\end{equation*}

la derivata prima è

\begin{equation*}
    S'_i(T)
    =
    b_i+2c_iT+3e_iT^2.
\end{equation*}

Le condizioni di continuità della derivata prima sono quindi
equazioni lineari nei coefficienti della spline.

I nodi interni sono $n-2$ e otteniamo pertanto altre

\begin{equation*}
    n-2
\end{equation*}

equazioni.

Con 15 nodi:

\begin{equation*}
    15-2=13.
\end{equation*}

Abbiamo quindi altre 13 equazioni.


\subsubsection{Terza condizione: continuità della derivata seconda}

Richiediamo una condizione ulteriore di regolarità.

Non vogliamo soltanto che i polinomi abbiano la stessa pendenza nei
punti nei quali si incontrano, ma anche che la curvatura della
funzione cambi in maniera regolare.

Imponiamo quindi la continuità della derivata seconda:

\begin{equation*}
    S''_{i-1}(T_i)=S''_i(T_i),
    \qquad i=2,\ldots,n-1.
    \label{eq:spline-second-derivative}
\end{equation*}

Poiché

\begin{equation*}
    S'_i(T)=b_i+2c_iT+3e_iT^2,
\end{equation*}

abbiamo

\begin{equation*}
    S''_i(T)=2c_i+6e_iT.
\end{equation*}

Anche queste condizioni sono lineari rispetto ai coefficienti
incogniti.

Poiché abbiamo $n-2$ nodi interni, otteniamo altre

\begin{equation*}
    n-2
\end{equation*}

equazioni.

Nel nostro esempio:

\begin{equation*}
    15-2=13.
\end{equation*}

Abbiamo quindi ottenuto altre 13 equazioni.


\subsubsection{Il conteggio delle equazioni}

A questo punto è utile verificare quante equazioni abbiamo costruito.

Le condizioni di interpolazione dei nodi forniscono

\begin{equation*}
    2(n-1)
\end{equation*}

equazioni.

La continuità della derivata prima fornisce

\begin{equation*}
    n-2
\end{equation*}

equazioni, mentre la continuità della derivata seconda ne fornisce
altre

\begin{equation*}
    n-2.
\end{equation*}

Complessivamente abbiamo quindi

\begin{equation*}
\begin{split}
2(n-1)+(n-2)+(n-2)
&=2n-2+n-2+n-2\\
&=4n-6.
\end{split}
\end{equation*}

Le incognite sono però

\begin{equation*}
    4(n-1)=4n-4.
\end{equation*}

Abbiamo dunque due incognite in più rispetto al numero di equazioni:

\begin{equation*}
    (4n-4)-(4n-6)=2.
\end{equation*}

Per determinare univocamente tutti i coefficienti della spline sono
quindi necessarie altre due condizioni.


\subsubsection{Condizioni al contorno e spline cubica naturale}

Le ultime due equazioni vengono ottenute imponendo delle condizioni
agli estremi dell'intervallo.

Le condizioni possibili sono diverse e vengono riassunte nella tabella \ref{tab:spline2cond}.

\begin{table}[htbp]
\centering
\caption{Confronto delle condizioni agli estremi delle spline cubiche}
\label{tab:condizioni_spline}
\renewcommand{\arraystretch}{1.5}
\begin{tabular}{ll}
\toprule
\textbf{Tipo di spline} & \textbf{Condizioni aggiuntive} \\
\midrule
Naturale & $S''(x_0)=0,\quad S''(x_n)=0$ \\
Vincolata (clamped) & $S'(x_0)=a,\quad S'(x_n)=b$ \\
Not-a-knot & $S'''(x_1^-)=S'''(x_1^+),\quad S'''(x_{n-1}^-)=S'''(x_{n-1}^+)$ \\
Derivata seconda assegnata & $S''(x_0)=a,\quad S''(x_n)=b$ \\
Periodica & $S'(x_0)=S'(x_n),\quad S''(x_0)=S''(x_n)$ \\
\bottomrule
\end{tabular}
\label{tab:spline2cond}
\end{table}


Una possibile scelta è quella della \emph{natural cubic spline}, o
spline cubica naturale.

In questo caso imponiamo che la derivata seconda sia nulla nei due
estremi:

\begin{equation*}
    S''_1(T_1)=0
\end{equation*}

e

\begin{equation*}
    S''_{n-1}(T_n)=0.
\end{equation*}

Queste due condizioni completano il sistema.

Abbiamo infatti

\begin{equation*}
    4n-6+2=4n-4
\end{equation*}

equazioni, esattamente pari al numero

\begin{equation*}
    4(n-1)=4n-4
\end{equation*}

di coefficienti incogniti.


\subsubsection{La spline come sistema lineare}

Possiamo ora tornare al nostro esempio con 15 BOT.

Il numero totale delle incognite è

\begin{equation*}
    4(15-1)=56.
\end{equation*}

Le diverse condizioni producono:

\begin{center}
\begin{tabular}{lr}
\hline
Condizione & Numero di equazioni\\
\hline
Interpolazione dei nodi & 28\\
Continuità di $S'$ & 13\\
Continuità di $S''$ & 13\\
Condizioni al contorno & 2\\
\hline
Totale & 56\\
\hline
\end{tabular}
\end{center}

Abbiamo quindi

\begin{equation*}
    56\text{ equazioni}
    \qquad
    56\text{ incognite}
\end{equation*}

Se raccogliamo tutti i coefficienti dei polinomi in un unico vettore
$x$, tutte le condizioni precedenti possono essere organizzate in
forma matriciale.

Il problema assume quindi la forma

\begin{equation*}
    A_{56\times56}x_{56\times1}
    =
    b_{56\times1}
\end{equation*}

Anche una procedura apparentemente complessa come la costruzione di
una spline cubica viene quindi ricondotta allo stesso problema
computazionale incontrato più volte nel corso:

\begin{equation*}
    Ax=b.
\end{equation*}

Questo risultato è particolarmente importante dal punto di vista
didattico. Il problema finanziario ci ha condotto alla necessità di
interpolare una funzione di sconto; la necessità di utilizzare molti
punti ci ha spinto a considerare le spline; la costruzione della
spline, a sua volta, si traduce nuovamente nella soluzione di un
sistema di equazioni lineari.

Le tecniche numeriche studiate in precedenza non costituiscono quindi
strumenti isolati, ma possono essere utilizzate come componenti di
procedure computazionali più articolate.

\subsubsection{Applicazione ai dati dei BOT}

Figure \ref{fig:spline} shows the cubic spline interpolation for the BOT example.

\begin{figure}
\centering
\includegraphics[scale=0.4]{fig_interp_spline.pdf}
\caption{cubic spline interpolation}
\label{fig:spline}
\end{figure}

\cleardoublepage

\section{Costruzione della struttura a termine con i BTP}
\label{sec:term-structure-btp}

Nelle sezioni precedenti abbiamo introdotto gli strumenti necessari
per affrontare il problema della costruzione della struttura a termine
dei tassi di interesse.

I BOT ci hanno permesso di osservare direttamente alcuni punti della
funzione di sconto. Abbiamo poi introdotto l'interpolazione per
ottenere valori della funzione in corrispondenza di scadenze non
direttamente osservate. In particolare, le spline cubiche permettono
di rappresentare la funzione di sconto mediante polinomi di basso
grado definiti localmente.

Possiamo ora tornare ai BTP e utilizzare questi strumenti per
costruire un'applicazione più vicina ai dati finanziari reali.

L'obiettivo è utilizzare congiuntamente:

\begin{itemize}
    \item i prezzi osservati dei BTP;
    \item i flussi di cassa generati dai titoli;
    \item la rappresentazione mediante spline della funzione di sconto;
\end{itemize}

per ricavare informazioni sulla struttura a termine.


\subsection{L'equazione di pricing di un BTP}

Consideriamo un BTP che genera $m$ pagamenti futuri alle date

\begin{equation*}
    T_1,T_2,\ldots,T_m.
\end{equation*}

Indicando con $CF_j$ il flusso di cassa alla data $T_j$, il prezzo
del titolo può essere scritto come

\begin{equation*}
    P
    =
    \sum_{j=1}^{m}
    CF_jd(T_j),
    \label{eq:btp-pricing-general}
\end{equation*}

dove $d(T_j)$ rappresenta il fattore di sconto corrispondente alla
data del pagamento.

Per un BTP nominale a tasso fisso, i pagamenti intermedi sono
costituiti dalle cedole, mentre l'ultimo pagamento comprende anche
il rimborso del valore nominale.

Se indichiamo con $C$ la cedola e assumiamo un valore nominale pari
a 100, possiamo quindi scrivere

\begin{equation*}
    P
    =
    Cd(T_1)
    +Cd(T_2)
    +\cdots+
    (100+C)d(T_m).
    \label{eq:btp-pricing-fixed}
\end{equation*}

Il problema è che i valori

\begin{equation*}
    d(T_1),d(T_2),\ldots,d(T_m)
\end{equation*}

non sono direttamente osservabili per tutte le date di pagamento.


\subsection{Interpolazione polinomiale}

Consideriamo ora un'applicazione nella quale utilizziamo i prezzi di mercato dei BTP per ricostruire la funzione di sconto mediante un polinomio. A differenza dell'interpolazione mediante spline cubiche, in questo caso rappresentiamo la funzione attraverso un unico polinomio definito sull'intero intervallo delle scadenze considerate.

Lo script R utilizza quattro BTP con scadenza il 1\textdegree{} aprile degli anni 2027, 2028, 2030 e 2031. Per ciascun titolo vengono estratti dal dataset il prezzo di mercato, la data di scadenza e il tasso cedolare annuo. Quest'ultimo viene ricavato automaticamente dalla descrizione del titolo e diviso per due per ottenere la cedola semestrale, assumendo un valore nominale pari a 100.

\paragraph{Preparazione dei prezzi e dei flussi di cassa.}
I prezzi osservati sul mercato sono generalmente espressi come \emph{clean prices}, ossia al netto degli interessi maturati dall'ultima data di pagamento della cedola. Per valutare i titoli mediante l'attualizzazione dei flussi di cassa futuri, dobbiamo invece utilizzare i \emph{dirty prices}, che comprendono anche il rateo di interesse maturato.

Lo script calcola quindi, per ciascun BTP,
\[
P_i^{\mathrm{dirty}}=P_i^{\mathrm{clean}}+R_i,
\]
dove $R_i$ rappresenta il rateo di interesse, ottenuto moltiplicando la cedola semestrale per la frazione del periodo cedolare già trascorsa.

Successivamente vengono ricostruite le date dei pagamenti futuri. Partendo dalla data di scadenza di ciascun titolo, lo script procede all'indietro con intervalli di sei mesi, individuando tutte le date successive alla data di valutazione.

Per ciascuna data viene calcolato il tempo residuo alla scadenza, espresso in anni:
\[
T_{ij}=\frac{d_{ij}-d_0}{365},
\]
dove $d_{ij}$ è la data del pagamento $j$ del titolo $i$, mentre $d_0$ è la data di valutazione.

I flussi di cassa sono costituiti dalle cedole semestrali e, in corrispondenza dell'ultima scadenza, dal rimborso del capitale.

\paragraph{Rappresentazione polinomiale della funzione di sconto.}
Supponiamo di rappresentare la funzione di sconto mediante un polinomio di terzo grado:
\[
d(T)=a_0+a_1T+a_2T^2+a_3T^3.
\]
I quattro coefficienti $a_0,a_1,a_2,a_3$ sono incogniti e devono essere determinati utilizzando i prezzi osservati dei BTP.

Ricordiamo che il prezzo di ciascun titolo è uguale alla somma dei valori attuali dei suoi flussi di cassa:
\[
P_i=\sum_{j=1}^{m_i}CF_{ij}d(T_{ij}).
\]
Sostituendo la rappresentazione polinomiale otteniamo:
\[
P_i=\sum_{j=1}^{m_i}CF_{ij}
\left(a_0+a_1T_{ij}+a_2T_{ij}^2+a_3T_{ij}^3\right).
\]
Raccogliendo i coefficienti incogniti:
\begin{align*}
P_i={}&a_0\sum_{j=1}^{m_i}CF_{ij}
+a_1\sum_{j=1}^{m_i}CF_{ij}T_{ij}\\
&+a_2\sum_{j=1}^{m_i}CF_{ij}T_{ij}^2
+a_3\sum_{j=1}^{m_i}CF_{ij}T_{ij}^3.
\end{align*}
L'equazione è lineare nei quattro coefficienti del polinomio.

\paragraph{Costruzione e soluzione del sistema lineare.}
Lo script costruisce una matrice $B$ nella quale ogni riga corrisponde a un BTP e ogni colonna contiene la somma dei flussi di cassa moltiplicati per una determinata potenza del tempo:
\[
B_i=\begin{pmatrix}
\displaystyle\sum_j CF_{ij} &
\displaystyle\sum_j CF_{ij}T_{ij} &
\displaystyle\sum_j CF_{ij}T_{ij}^2 &
\displaystyle\sum_j CF_{ij}T_{ij}^3
\end{pmatrix}.
\]
Poiché utilizziamo quattro BTP e quattro coefficienti incogniti, otteniamo una matrice quadrata $B$ di dimensione $4\times4$.

Definendo
\[
\mathbf a=\begin{pmatrix}a_0\\a_1\\a_2\\a_3\end{pmatrix},
\qquad
\mathbf P=\begin{pmatrix}P_1\\P_2\\P_3\\P_4\end{pmatrix},
\]
il problema può essere scritto nella forma:
\[
\boxed{B\mathbf a=\mathbf P}.
\]
Lo script risolve questo sistema mediante la funzione \texttt{solve()} di R. Se la matrice $B$ è invertibile, la soluzione è unica e i coefficienti ottenuti permettono di riprodurre esattamente i quattro prezzi utilizzati nella calibrazione.

\paragraph{Ricostruzione della curva dei fattori di sconto.}
Una volta determinati i coefficienti, possiamo ricostruire la funzione:
\[
\hat d(T)=\hat a_0+\hat a_1T+\hat a_2T^2+\hat a_3T^3.
\]
Lo script calcola i valori della funzione su una griglia di 500 scadenze comprese tra zero e la scadenza massima dei titoli selezionati, rappresentando graficamente la curva ottenuta.

Il risultato è una funzione continua che permette di calcolare un fattore di sconto anche in corrispondenza di scadenze diverse da quelle dei BTP utilizzati. La figura \ref{fig:btp_poli} mostra la funzione.

\begin{figure}
\centering
\includegraphics[scale=0.4]{fig_interp_btp_poli.pdf}
\caption{interpolazione polinomiale con BTP}
\label{fig:btp_poli}
\end{figure}




È importante osservare che, in questa applicazione, \textbf{non vengono imposte condizioni di raccordo tra polinomi diversi}, poiché la funzione di sconto è rappresentata da un unico polinomio cubico globale.

Inoltre, il procedimento non impone esplicitamente la condizione $d(0)=1$, né garantisce che i fattori di sconto siano positivi e decrescenti. Queste proprietà devono quindi essere verificate sulla curva ottenuta.

L'esempio mostra come, partendo dai prezzi e dai flussi di cassa di un insieme di obbligazioni, sia possibile ricondurre la costruzione della funzione di sconto alla soluzione di un sistema lineare. Rappresenta inoltre un utile termine di confronto con l'approccio basato sulle spline cubiche, nel quale la funzione viene descritta mediante più polinomi locali collegati da opportune condizioni di raccordo.

\subsection{Interpolazione mediante spline}

Nella sezione precedente abbiamo costruito una spline cubica per
rappresentare la funzione di sconto.

Supponiamo che la data $T_j$ appartenga all'intervallo compreso tra
due nodi consecutivi:

\begin{equation*}
    T_i\leq T_j\leq T_{i+1}.
\end{equation*}

All'interno di questo intervallo il fattore di sconto viene
rappresentato dal polinomio

\begin{equation*}
    d(T)
    =
    a_i+b_iT+c_iT^2+e_iT^3.
    \label{eq:discount-spline-segment}
\end{equation*}

Possiamo quindi calcolare il fattore di sconto relativo alla data
$T_j$ come

\begin{equation*}
    d(T_j)
    =
    a_i+b_iT_j+c_iT_j^2+e_iT_j^3.
    \label{eq:discount-at-cashflow}
\end{equation*}

Una volta conosciuti i coefficienti della spline, possiamo pertanto
ottenere il fattore di sconto corrispondente a qualsiasi data
compresa nell'intervallo coperto dai dati.


\subsubsection{Dalla spline all'equazione di pricing}

Consideriamo nuovamente l'equazione

\begin{equation*}
    P
    =
    \sum_{j=1}^{m}
    CF_jd(T_j).
\end{equation*}

Per ciascuna data $T_j$, il valore $d(T_j)$ è determinato dal tratto
della spline che contiene quella data.

Se, ad esempio, $T_j$ appartiene all'intervallo
$[T_i,T_{i+1}]$, possiamo sostituire

\begin{equation*}
    d(T_j)
    =
    a_i+b_iT_j+c_iT_j^2+e_iT_j^3.
\end{equation*}

Il contributo del flusso di cassa $CF_j$ al prezzo diventa quindi

\begin{equation*}
    CF_jd(T_j)
    =
    CF_j
    \left(
    a_i+b_iT_j+c_iT_j^2+e_iT_j^3
    \right).
\end{equation*}

Sviluppando il prodotto otteniamo

\begin{equation*}
    CF_jd(T_j)
    =
    CF_ja_i
    +
    CF_jT_jb_i
    +
    CF_jT_j^2c_i
    +
    CF_jT_j^3e_i.
    \label{eq:cf-spline-linear}
\end{equation*}

Questa espressione mette in evidenza una proprietà fondamentale:
l'equazione è \emph{lineare nei coefficienti della spline}.

Le quantità

\begin{equation*}
    CF_j,\qquad
    CF_jT_j,\qquad
    CF_jT_j^2,\qquad
    CF_jT_j^3
\end{equation*}

sono infatti note. Le incognite sono invece

\begin{equation*}
    a_i,\qquad b_i,\qquad c_i,\qquad e_i.
\end{equation*}

Questo permette di trasformare l'equazione di pricing di ciascun
BTP in un'equazione lineare.


\subsubsection{Il contributo dei diversi intervalli}

Un BTP genera normalmente più pagamenti e le relative date possono
appartenere a intervalli differenti della spline.

Consideriamo schematicamente tre intervalli:

\begin{equation*}
[T_1,T_2],
\qquad
[T_2,T_3],
\qquad
[T_3,T_4].
\end{equation*}

A essi corrispondono tre polinomi:

\begin{align}
S_1(T)&=a_1+b_1T+c_1T^2+e_1T^3,\\
S_2(T)&=a_2+b_2T+c_2T^2+e_2T^3,\\
S_3(T)&=a_3+b_3T+c_3T^2+e_3T^3.
\end{align}

Se un pagamento avviene nel primo intervallo, utilizzeremo i
coefficienti di $S_1$. Se avviene nel secondo intervallo,
utilizzeremo quelli di $S_2$, e così via.

L'equazione di pricing del titolo può quindi coinvolgere
coefficienti appartenenti a diversi segmenti della spline.

Possiamo raccogliere tutti questi coefficienti in un unico vettore:

\begin{equation*}
a=
\begin{pmatrix}
a_1\\
b_1\\
c_1\\
e_1\\
a_2\\
b_2\\
c_2\\
e_2\\
\vdots
\end{pmatrix}.
\label{eq:spline-coefficient-vector-finance}
\end{equation*}

Per ogni BTP possiamo allora costruire un vettore riga che contiene
i coefficienti noti che moltiplicano gli elementi di $a$.


\subsubsection{Un'equazione lineare per ciascun BTP}

Indichiamo con $B_i$ il vettore riga costruito utilizzando i flussi
di cassa e le date di pagamento del BTP $i$.

L'equazione di pricing può essere scritta come

\begin{equation*}
    B_i a=P_i,
    \label{eq:single-btp-linear}
\end{equation*}

dove:

\begin{itemize}
    \item $P_i$ è il prezzo osservato del titolo;
    \item $B_i$ contiene informazioni sui cash flow e sulle
          rispettive date;
    \item $a$ contiene i coefficienti incogniti della
          rappresentazione della funzione di sconto.
\end{itemize}

Ogni titolo osservato fornisce quindi una nuova equazione lineare.

Considerando simultaneamente più BTP possiamo raccogliere le
equazioni in un sistema:

\begin{equation*}
\begin{pmatrix}
B_1\\
B_2\\
\vdots\\
B_N
\end{pmatrix}
a
=
\begin{pmatrix}
P_1\\
P_2\\
\vdots\\
P_N
\end{pmatrix}.
\end{equation*}

Definendo

\begin{equation*}
B=
\begin{pmatrix}
B_1\\
B_2\\
\vdots\\
B_N
\end{pmatrix}
\end{equation*}

e

\begin{equation*}
P=
\begin{pmatrix}
P_1\\
P_2\\
\vdots\\
P_N
\end{pmatrix},
\end{equation*}

otteniamo

\begin{equation*}
    Ba=P
    \label{eq:BaP}
\end{equation*}

Ancora una volta, il problema finanziario è stato trasformato nella
forma generale

\begin{equation*}
    Ax=b.
\end{equation*}


\subsubsection{Interpretazione della matrice $B$}

È utile soffermarsi sul significato della matrice $B$.

Le righe della matrice corrispondono ai titoli utilizzati
nell'applicazione. Le colonne corrispondono invece ai coefficienti
della funzione utilizzata per rappresentare i fattori di sconto.

Gli elementi di $B$ vengono costruiti combinando:

\begin{equation*}
\text{cash flow}
\quad+\quad
\text{date dei cash flow}
\quad+\quad
\text{forma della spline}.
\end{equation*}

La matrice $B$ rappresenta quindi il collegamento tra le
caratteristiche osservabili dei titoli e i parametri della funzione
di sconto che vogliamo determinare.

Questo passaggio è importante perché mostra come la costruzione di
un sistema lineare non sia un'operazione puramente formale.
La matrice dei coefficienti deve essere costruita a partire dalla
struttura economica o finanziaria del problema.



\subsubsection{Le condizioni di raccordo e il completamento del sistema lineare}

Nelle sezioni precedenti abbiamo mostrato come le equazioni di pricing dei BTP possano essere trasformate in equazioni lineari nei coefficienti della funzione di sconto. Tuttavia, quando rappresentiamo tale funzione mediante una spline cubica, dobbiamo considerare anche le condizioni che collegano i polinomi definiti nei diversi intervalli.

Supponiamo di suddividere l'orizzonte temporale in $n$ intervalli. In ciascun intervallo la funzione di sconto è rappresentata da un polinomio cubico:
\[
S_i(T)=a_i+b_iT+c_iT^2+e_iT^3,
\qquad i=1,\ldots,n.
\]
Poiché ogni polinomio contiene quattro coefficienti, il numero complessivo delle incognite è $4n$.

Per ottenere una spline cubica, i polinomi adiacenti devono raccordarsi in modo regolare. In particolare, in ogni nodo interno $T_i$ imponiamo la continuità della funzione e delle sue derivate prima e seconda:
\[
\begin{aligned}
S_i(T_i)&=S_{i+1}(T_i),\\
S_i'(T_i)&=S_{i+1}'(T_i),\\
S_i''(T_i)&=S_{i+1}''(T_i).
\end{aligned}
\]
Queste sono le cosiddette \emph{condizioni di raccordo} (\emph{smooth pasting conditions}). Poiché i nodi interni sono $n-1$, otteniamo complessivamente $3(n-1)$ equazioni lineari.

Le condizioni di raccordo non sono però sufficienti, da sole, a determinare tutti i coefficienti. Infatti, rimangono
\[
4n-3(n-1)=n+3
\]
gradi di libertà, che devono essere determinati utilizzando le equazioni di pricing dei BTP e, se necessario, ulteriori condizioni.

Ad esempio, consideriamo una spline composta da tre polinomi cubici e supponiamo di disporre dei prezzi di quattro BTP. Abbiamo dodici coefficienti incogniti. Le quattro equazioni di pricing, insieme alle sei condizioni di raccordo nei due nodi interni, forniscono dieci equazioni. Per completare il sistema possiamo imporre altre due condizioni, ad esempio:
\[
d(0)=1,\qquad d''(T_{\max})=0.
\]
La prima esprime una proprietà fondamentale dei fattori di sconto: un pagamento immediato non deve essere scontato. La seconda impone che la derivata seconda della funzione di sconto sia nulla all'estremità destra dell'intervallo. Quest'ultima è una possibile condizione al contorno, non una proprietà finanziaria necessaria.

Otteniamo così un sistema di dodici equazioni in dodici incognite.

In generale, possiamo raccogliere le equazioni di pricing e le condizioni aggiuntive in un unico sistema lineare:
\[
\boxed{
\begin{pmatrix}
B\\ C
\end{pmatrix}
\mathbf a=
\begin{pmatrix}
\mathbf P\\ \mathbf c
\end{pmatrix}
}
\]
dove $B$ è la matrice costruita a partire dai flussi di cassa dei BTP, $C$ contiene i coefficienti delle condizioni di raccordo e delle condizioni al contorno, mentre $\mathbf c$ raccoglie i termini noti di queste ultime.

Quando il numero di equazioni coincide con quello delle incognite e la matrice completa è invertibile, il sistema ammette una soluzione unica. Se invece disponiamo di un numero maggiore di equazioni di pricing, possiamo ricorrere a metodi di stima ai minimi quadrati, mantenendo i vincoli necessari per garantire la regolarità della spline.

In questo modo, la costruzione della struttura a termine si riconduce a un problema di algebra lineare nel quale le equazioni di pricing incorporano le informazioni finanziarie, mentre le condizioni di raccordo assicurano la regolarità matematica della funzione di sconto.




\subsubsection{Dalla soluzione del sistema alla funzione di sconto}

Una volta costruiti $B$ e $P$, possiamo risolvere

\begin{equation*}
    Ba=P
\end{equation*}

utilizzando le tecniche numeriche studiate nelle sezioni precedenti.

La soluzione fornisce il vettore dei coefficienti

\begin{equation*}
    a^*.
\end{equation*}

Conoscendo questi coefficienti possiamo ricostruire i diversi
polinomi della spline e quindi la funzione

\begin{equation*}
    d(T).
\end{equation*}

Il procedimento può essere sintetizzato come

\begin{equation*}
\text{prezzi e cash flow}
\longrightarrow
Ba=P
\longrightarrow
a^*
\longrightarrow
d(T).
\end{equation*}

A differenza del semplice esempio iniziale, non stiamo quindi
stimando separatamente un fattore di sconto per ciascuna data.
Determiniamo invece i parametri di una funzione che permette di
calcolare il fattore di sconto per le diverse scadenze.

I risultati dell'interpolazione con tre intervalli sono riportati nella tabella \ref{tab:coefficienti_cubici}.

\begin{table}[htbp]
\centering
\caption{Coefficienti dei polinomi cubici}
\label{tab:coefficienti_cubici}
\begin{tabular}{rrrrrr}
\toprule
Intervallo & $[T_i,T_{i+1}]$ & $a_i$ & $b_i$ & $c_i$ & $e_i$ \\
\midrule
1 & [0,1] &1.0000000 & -0.02591576 & $-9.123480 \times 10^{-3}$ & $2.956755 \times 10^{-3}$ \\
2 & [1,3] &1.0030381 & -0.03502999 & $-9.252434 \times 10^{-6}$ & $-8.132081 \times 10^{-5}$ \\
3 & [3,5] &0.9975073 & -0.02949920 & $-1.852849 \times 10^{-3}$ & $1.235233 \times 10^{-4}$ \\
\bottomrule
\end{tabular}
\end{table}

La figura \ref{fig:btp_spline} mostra la linea interpolante.

\begin{figure}
\centering
\includegraphics[scale=0.4]{fig_interp_btp_spline.pdf}
\caption{interpolazione spline con BTP}
\label{fig:btp_spline}
\end{figure}



\subsubsection{Dalla funzione di sconto alla curva spot}

Una volta ricostruita la funzione $d(T)$ possiamo trasformare i
fattori di sconto nei corrispondenti tassi spot.

Ricordiamo che

\begin{equation*}
    d(T)
    =
    \frac{1}{(1+r(T))^T}.
\end{equation*}

Risolvendo rispetto al tasso otteniamo

\begin{equation*}
    r(T)
    =
    d(T)^{-1/T}-1
    \label{eq:spline-discount-to-spot}
\end{equation*}

Possiamo quindi scegliere una griglia di maturities

\begin{equation*}
    T_1^*,T_2^*,\ldots,T_K^*
\end{equation*}

e, per ciascuna di esse:

\begin{enumerate}
    \item calcolare il fattore di sconto $d(T_k^*)$;
    \item trasformarlo nel corrispondente tasso spot
          $r(T_k^*)$.
\end{enumerate}

Le coppie

\begin{equation*}
    \left(T_k^*,r(T_k^*)\right)
\end{equation*}

permettono infine di rappresentare graficamente la struttura a
termine dei tassi di interesse.


\subsection{Il ruolo dei BOT e dei BTP}

Possiamo ora comprendere meglio il ruolo svolto dai due tipi di
titoli utilizzati nell'applicazione.

I BOT, essendo titoli zero-coupon, permettono di ricavare
direttamente alcuni punti della funzione di sconto:

\begin{equation*}
    d(T)=\frac{P_T}{100}.
\end{equation*}

I BTP forniscono invece informazioni aggiuntive attraverso
l'intera sequenza dei loro flussi di cassa:

\begin{equation*}
    P_i
    =
    \sum_j CF_{ij}d(T_{ij}).
\end{equation*}

La funzione di interpolazione permette di collegare queste
informazioni e di rappresentare $d(T)$ anche in corrispondenza
delle date per le quali non osserviamo direttamente un titolo
zero-coupon.

La costruzione della curva può quindi essere vista come
l'integrazione di tre elementi:

\begin{equation*}
\begin{array}{ccc}
\text{BOT}
&
+
&
\text{BTP}\\[1mm]
\downarrow&&\downarrow\\
\text{punti della funzione di sconto}
&&
\text{equazioni di pricing}\\[2mm]
&\searrow\quad\swarrow&\\[-1mm]
&d(T)&
\end{array}
\end{equation*}


\subsection{Dal problema finanziario al problema computazionale}

L'intero percorso seguito nell'applicazione può essere ora
ricostruito.

Siamo partiti da una domanda finanziaria:

\begin{quote}
come possiamo ricavare la struttura a termine dei tassi di interesse
utilizzando i prezzi dei titoli osservati sul mercato?
\end{quote}

Abbiamo inizialmente considerato il caso semplice dei BOT, per i
quali il prezzo permette di ottenere direttamente un fattore di
sconto.

Successivamente abbiamo incontrato un problema: le scadenze
disponibili sul mercato non coincidono necessariamente con tutte le
date nelle quali abbiamo bisogno di conoscere il fattore di sconto.

Questo ci ha condotto all'interpolazione:

\begin{equation*}
\text{punti osservati}
\longrightarrow
\text{funzione }d(T).
\end{equation*}

L'interpolazione polinomiale globale ha evidenziato possibili
problemi di condizionamento. Abbiamo quindi introdotto le spline,
che rappresentano la funzione mediante polinomi di basso grado
definiti localmente.

Infine, siamo tornati alle equazioni di pricing dei BTP. Sostituendo
in tali equazioni la rappresentazione scelta per $d(T)$ abbiamo
ottenuto un sistema lineare:

\begin{equation*}
    Ba=P.
\end{equation*}

La soluzione del sistema fornisce i coefficienti necessari per
ricostruire la funzione di sconto e, successivamente, la curva dei
tassi spot.


\subsection{Il messaggio computazionale}

L'applicazione finanziaria può essere sintetizzata attraverso la
sequenza

\begin{equation*}
\text{prezzi + cash flow}
\longrightarrow
Ba=P
\longrightarrow
a
\longrightarrow
d(T)
\longrightarrow
r(T).
\end{equation*}

Il passaggio centrale è ancora una volta la soluzione di un sistema
lineare.

Tuttavia, questa applicazione mette in evidenza anche un secondo
aspetto importante del calcolo numerico: la qualità della soluzione
non dipende soltanto dall'algoritmo utilizzato per risolvere il
sistema.

Dipende anche da \emph{come il problema viene rappresentato}.

Nel nostro caso, la scelta della funzione utilizzata per descrivere
$d(T)$ determina la struttura del sistema lineare che dobbiamo
risolvere e influenza le proprietà della curva ottenuta.

L'interpolazione polinomiale e le spline non sono quindi semplici
passaggi accessori, ma fanno parte della formulazione numerica del
problema.

L'applicazione finanziaria riunisce così i principali concetti
introdotti in questa parte del corso:

\begin{equation*}
\boxed{
\begin{array}{c}
\text{problema finanziario}\\
\downarrow\\
\text{rappresentazione matematica}\\
\downarrow\\
\text{interpolazione}\\
\downarrow\\
\text{sistema lineare}\\
\downarrow\\
\text{soluzione numerica}\\
\downarrow\\
\text{interpretazione finanziaria}
\end{array}}
\end{equation*}

Il risultato finale non è quindi soltanto una curva dei tassi.
L'esempio mostra, più in generale, come un problema finanziario reale
possa essere trasformato progressivamente in una sequenza di problemi
matematici che il computer è in grado di risolvere.

